\documentclass[aspectratio=169,11pt]{beamer}

% Compile from this directory: latexmk -pdf -outdir=build presentation.tex
\usepackage[T1]{fontenc}
\usepackage{lmodern}
\usepackage{amsmath,amssymb,mathtools}
\usepackage{tikz-cd}

\usetheme{Madrid}
\usefonttheme{professionalfonts}
\definecolor{accent}{HTML}{176B70}
\setbeamercolor{structure}{fg=accent}
\setbeamercolor{alerted text}{fg=accent}
\setbeamertemplate{navigation symbols}{}
\setbeamertemplate{footline}[frame number]
\setbeamertemplate{bibliography item}[text]

% Metadata: replace the bracketed fields before presenting.
\title[Rigid Local Systems]{Rigid Local Systems}
\subtitle{Candidacy Presentation}
\author{[Your name]}
\institute{[Department] \\ {[University]}}
\date{[Presentation date]}

% Notation: keep consistent with writeup/writeup.tex.
\newcommand{\C}{\mathbb{C}}
\newcommand{\Z}{\mathbb{Z}}
\newcommand{\PP}{\mathbb{P}}
\DeclareMathOperator{\GL}{GL}
\DeclareMathOperator{\End}{End}
\DeclareMathOperator{\Hom}{Hom}
\DeclareMathOperator{\rank}{rank}
\DeclareMathOperator{\rig}{rig}
\newcommand{\placeholder}[1]{{\color{gray}\textit{[#1]}}}

\begin{document}

\begin{frame}
  \titlepage
\end{frame}

\begin{frame}{Roadmap}
  \tableofcontents
\end{frame}

\section{Motivation}

\begin{frame}{The guiding question}
  \begin{block}{Question}
    When does local monodromy determine a local system up to isomorphism?
  \end{block}
  \medskip
  \placeholder{Explain the motivation and give one concrete example.}
\end{frame}

\section{Background and definitions}

\begin{frame}{Geometric setup}
  Let
  \[
    U = \PP^1(\C) \setminus S,
    \qquad S = \{s_1,\ldots,s_m\}.
  \]
  \begin{itemize}
    \item \placeholder{Specify the coefficient field and rank.}
    \item \placeholder{Define the local systems under consideration.}
    \item \placeholder{Explain the local monodromy at each puncture.}
  \end{itemize}
\end{frame}

\begin{frame}{Physical rigidity}
  \begin{definition}[Physical rigidity]
    \placeholder{Insert the precise definition, including the class of
    competing local systems and the meaning of the same local monodromy.}
  \end{definition}
  \medskip
  \placeholder{Distinguish local conjugacy from a single global conjugacy.}
\end{frame}

\section{Main result}

\begin{frame}{The main theorem}
  \begin{theorem}[\placeholder{Attribution and theorem number}]
    \placeholder{State every hypothesis and the exact conclusion.}
  \end{theorem}
  \medskip
  \placeholder{Explain the significance in one sentence.}
\end{frame}

\begin{frame}{Proof strategy}
  \begin{enumerate}
    \item \placeholder{First reduction or key construction.}
    \item \placeholder{Central calculation or lemma.}
    \item \placeholder{Deduce the theorem.}
  \end{enumerate}
\end{frame}

\begin{frame}{A key step}
  \begin{lemma}
    \placeholder{State the key lemma.}
  \end{lemma}
  \begin{proof}[Proof idea]
    \placeholder{Give the main argument; move lengthy details to the appendix.}
  \end{proof}
\end{frame}

\section{Examples and next steps}

\begin{frame}{A worked example}
  \placeholder{Choose a low-rank example and compute its local data.}
  \[
    A_1\cdots A_m = I,
    \qquad A_i \in \GL_n(\C).
  \]
  \placeholder{Explain what the theorem tells us in this example.}
\end{frame}

\begin{frame}{Next steps}
  \begin{itemize}
    \item \placeholder{A mathematical question to investigate.}
    \item \placeholder{A formalization task, with precise remaining dependencies.}
    \item \placeholder{A concrete milestone.}
  \end{itemize}
\end{frame}

\begin{frame}{References}
  % Replace with bibliographic entries and use \cite{key} on relevant slides.
  \placeholder{Add the references used in the talk.}
\end{frame}

\appendix
\section*{Backup slides}

\begin{frame}{Backup: details and computations}
  \placeholder{Put additional proofs, examples, or definitions here.}
  % A frame containing tikzcd should use \begin{frame}[fragile]{Title}.
\end{frame}

\end{document}
